The concentration of blink information at Fp1 and Fp2 agrees with earlier driving studies that detected blinks from a single frontopolar or forehead channel and used them to estimate driver drowsiness or fatigue
~\citep{ko2020eyeblink,shahbakhti2022simultaneous,shahbakhti2023fusion}.
Those studies fixed the frontopolar channel by design. In the present benchmark, the same two electrodes came out on top in every configuration without any such constraint, and neither lost significant performance when the pipeline received only the Frontal region or that electrode alone. Significant losses were common elsewhere under the same reduction, including at the neighbouring frontal electrodes. Because the pipeline derives its candidate epochs and thresholds entirely from the electrodes supplied, the frontopolar advantage does not rely on information borrowed from other channels. The pattern held on both datasets, although the C3--C4 asymmetry was specific to Internal and may reflect a recording or montage characteristic of that dataset. These results also have implications for reduced-channel drowsiness systems, which commonly use occipital electrodes because alpha activity, a well-established marker of driver fatigue~\citep{simon2011eeg}, is prominent over the occipital region~\citep{lin2014wireless,li2015context}. Occipital electrodes performed far below the frontopolar electrodes in every configuration tested here, so a system relying on an occipital channel alone would capture alpha activity but miss most of the blink information available at frontopolar sites. That information is not redundant, since adding blink features to single-channel forehead EEG improved the prediction of driver fatigue in a simulated driving task~\citep{ko2020eyeblink}.
